\documentclass[exercice]{thedocument}%
\usepackage{texmaths-tikz}%
\startdatakeys
\source{}
\cycle{}
\competencesDuSocle{}
\connaissancesRequises{}
\competencesTravaillees{}
\enddatakeys
\begin{document}%
\begin{flushleft}
    Dans chaque cas, tracer la médiatrice du segment [AB].\par Cette médiatrice
    est-elle un axe de symétrie de la figure ?
\end{flushleft}\par\noindent
    \textbf{a)}\begin{tikzpicture}[baseline=(I.base),scale=1]
        \coordinate[label=above left:I](I)at(0,0);
        \coordinate[label=above right:J](J)at(4,0);
        \coordinate[label=below left:A](A)at(-0.5,-2);
        \coordinate[label=below right:B](B)at(3.5,-2);
        \filldraw[fill=Rainbow1-rouge](A)--(B)--(J)--(I)--cycle;
    \end{tikzpicture}\par\noindent
    \textbf{b)}\begin{tikzpicture}[rotate=20,baseline=(B.base),scale=1]
        \coordinate[label=left:A](I)at(0,0);
        \coordinate[label=above:D](J)at(4,0);
        \coordinate[label=below:C](A)at(0,-2);
        \coordinate[label=right:B](B)at(4,-2);
        \filldraw[fill=Rainbow1-jaune](A)--(B)--(J)--(I)--cycle;
    \end{tikzpicture}
\end{document}
